% Part: lambda-calculus
% Chapter: introduction
% Section: lambda-definability

\documentclass[../../../include/open-logic-section]{subfiles}

\begin{document}

\olfileid{lam}{int}{rep}
\olsection{Fungsi Aritmetika kang \usetoken{S}{lambda definable}}

Kepiye kalkulus lambda bisa dadi modhèl komputasi?
Wiwitane, malah durung cetha kepriye ukara iki
kudu dimangerteni. Kanggo ngrembug komputabilitas
ing wilangan asli, kita perlu nemokake representasi
kang trep kanggo wilangan-wilangan iku. Ing ngisor
iki salah siji representasi kang bisa digunakake
kanthi becik, luwih becik tinimbang kang bisa
dikira sadurunge.

\begin{defn}
Kanggo saben wilangan asli~$n$, tetepna
\emph{numeral Church} $\num{n}$ minangka
suku lambda $\lambd[x][\lambd[y][(x(x(x(\dots
x(y)))))]]$, kanthi cacah sakabehe ana $n$ aksara $x$.
\end{defn}

Suku-suku $\num{n}$ iku ``iterator'': yen
diwenehi input $f$, $\num{n}$ ngasilake
fungsi kang masangake $y$ karo $f^n(y)$.
Gatekna yen saben numeral iku wis normal.
Saiki kita bisa nerangake apa tegese sawijining
suku lambda ``ngitung'' fungsi ing wilangan asli.

\begin{defn}
Upamana $f(x_0, \dots, x_{k-1})$ sawijining fungsi
parsial kanthi $k$ argumen saka $\Nat$ lan
asil ing $\Nat$. Kita nyebut suku-$\lambd$~$F$
\emph{!!{lambda define}s}~$f$ yen lan mung yen
kanggo saben urutan wilangan asli $n_0$,
\dots,~$n_{k-1}$,
\[
F\, \num{n_0}\, \num{n_1} \dots \num{n_{k-1}} \red \num{f(n_0, n_1, \dots,
  n_{k-1})}
\]
yen $f(n_0, \dots, n_{k-1})$ nduweni nilai,
dene $F\, \num{n_0}\, \num{n_1}
\dots \num{n_{k-1}}$ ora nduweni wujud
normal yen fungsi iku ora nduweni nilai
ing input kasebut.
\end{defn}

\begin{thm}
\ollabel{thm:lambda-def}
Fungsi $f$ iku fungsi parsial kang bisa diitung
yen lan mung yen fungsi iku
!!{lambda defined} dening sawijining suku lambda.
\end{thm}

\begin{explain}
Teorema iki rada nggumunake. Minangka modhèl
komputasi, kalkulus lambda iku kalkulus kang
cukup prasaja: mung ana rong operasi, yaiku
abstraksi lambda lan aplikasi!{} Nanging saka
sarana kang winates iki, saben prosedur komputasi
bisa diwujudake.
\end{explain}

\end{document}
